\input{../tex-inputs/latex-korrekturansicht-vorspann}

\section[Arthur Schnitzler an Gustav Schwarzkopf, 2. 8. 1914]{L04351 Arthur Schnitzler an Gustav Schwarzkopf, 2. 8. 1914}
\nopagebreak\mylabel{L04351v}
\rehead{ }\normalsize\beginnumbering\briefempfaengerindex{Schwarzkopf, Gustav@\textsc{Schwarzkopf, Gustav}!zzzSchnitzler, Arthur@\emph{von Arthur Schnitzler}!1914-08-022@{2. 8. 1914}|(be}
\toendnotes[C]{\smallbreak\pagebreak[2]}
\correspDesc{Versand  durch Arthur Schnitzler am 2. 8. 1914 in Celerina
\newline{}Übermittlung  am 2. 8. 1914 in Celerina
\newline{}Erhalt  durch Gustav Schwarzkopf im Zeitraum [3. 8. 1914 – 7. 8. 1914?] in Wien}\toendnotes[C]{\smallbreak}
\Standort{DLA, A:Schnitzler, HS.1985.1.1897.}
\physDesc{Bildpostkarte, 421 Zeichen
\newline{}Handschrift: , deutsche Kurrent
\newline{}Versand: Stempel: »\nobreak{}\oindex{Celerina@\textbf{Celerina}|pwk}\textcolor{gray}{Celerina
                                          (Graubünden)}, 2. VIII. 14, 9\nobreak{}«.  }\pstart{}{\pb}\textcolor{pink}{Austria}\oindex{Österreich@\textbf{Österreich}|pw}{}\ledrightnote{\textcolor{pink}{Österreich}}\pend{}\pstart{}Hrn Gustav Schwarzkopf\pend{}\pstart{}\textcolor{pink}{Wien I}\oindex{I., Innere Stadt@\textbf{I., Innere Stadt}, \emph{Verwaltungsgebiet}|pw}{}\ledrightnote{\textcolor{pink}{I., Innere Stadt}}\pend{}\pstart{}\textcolor{pink}{Tiefer Graben 17}\oindex{Wien@\textbf{Wien}!I., Innere Stadt@\textbf{I., Innere Stadt}!Tiefer Graben 17@\textbf{Tiefer Graben 17}, \emph{Wohngebäude}|pw}{}\ledrightnote{\textcolor{pink}{Tiefer Graben 17}}\pend{}{\bigskip}
\pstart
           \noindent{}\centering{}{\pb}\textcolor{pink}{\textcolor{gray}{\textbf{Cresta Palace, Celerina}}}\oindex{Cresta Palace@\textbf{Cresta Palace}, \emph{Hotel}|pw}{}\ledrightnote{\textcolor{pink}{Cresta Palace}}\pend
           \vspace{1em}
\pstart
           \noindent{}{\pb}Sind Sie ſchon in \textcolor{pink}{Wien}\oindex{Wien@\textbf{Wien}, \emph{Verwaltungsgebiet}|pw}{}\ledrightnote{\textcolor{pink}{Wien}}? Oder noch? Wir ſind u bleiben (wie lange unbeſti{\geminationm}t) hier – wo es herrlich wäre, we{\geminationn} die Welt nicht durch Telegra{\geminationm}e verdüſtert und die Menſchheit toll geworden ſchiene
               –  Iſt Politik nicht eine phantaſtiſche {\pb}Sache? Beſonders
                  we{\geminationn} man denkt, daſs ſie am Tag drauf schon
               Weltgeſchichte iſt! Laſſen Sie was von ſich hören\pend
           
\pstart
           Wir grüßen herzlichſt{\\[\baselineskip]}Ihr \spacefill\mbox{A.}\pend
           \leftskip=0em{}
\pstart
           2. 8. 14\pend
           \selectlanguage{ngerman}\endnumbering\briefempfaengerindex{Schwarzkopf, Gustav@\textsc{Schwarzkopf, Gustav}!zzzSchnitzler, Arthur@\emph{von Arthur Schnitzler}!1914-08-022@{2. 8. 1914}|)be}\mylabel{L04351h}
\begin{anhang}
\end{anhang}\newcommand{\dateiname}{L04351}\newcommand{\titel}{Arthur Schnitzler an Gustav Schwarzkopf, 2. 8. 1914}\newcommand{\editorInnen}{Herausgegeben von Jahnke, SelmaMüller, Martin Anton}\input{../tex-inputs/latex-korrekturansicht-abspann}
